% TOP_PROBLEM: {"id": 3251, "title": "2-colouring a graph without a monochromatic maximum clique", "category_id": 3, "category": {"id": 3, "name": "graph_theory", "display_name": "Graph Theory", "description": "Problems involving graphs, networks, and their properties.", "slug": "graph-theory", "order_index": 3, "created_at": "2026-07-31T15:26:25.671Z"}, "status": "open", "problem_number": "OPG-56230", "set_id": 12, "statusQualification": "Non-empty is explicitly taken to mean at least one edge, avoiding the singleton-maximum-clique degeneracy."}
% FIRST_POSTED: 2026-10-02
% ATTEMPT_STATUS: unresolved
% UnsolvedMath ID 3251; source catalogue code OPG-56230
% !TeX program = lualatex
% Research attempt by GPT-6 Astra Ultra; no claim of a new complete solution.
\documentclass[11pt,a4paper]{article}
\usepackage[margin=25mm]{geometry}
\usepackage{fontspec}
\setmainfont{lmroman10-regular.otf}[BoldFont=lmroman10-bold.otf,
 ItalicFont=lmroman10-italic.otf,BoldItalicFont=lmroman10-bolditalic.otf]
\usepackage{amsmath,amssymb,mathtools}
\usepackage{unicode-math}
\setmathfont{latinmodern-math.otf}
\AtBeginDocument{\let\setminus\smallsetminus}
\usepackage{xurl,hyperref}
\hypersetup{colorlinks=true,urlcolor=blue}
\setcounter{secnumdepth}{0}
\setlength{\parindent}{0pt}
\setlength{\parskip}{5pt}
\setlength{\emergencystretch}{3em}
\begin{document}
\section*{2-colouring a graph without a monochromatic maximum clique}\label{3251}
{\small UnsolvedMath ID 3251 · OPG-56230 · Graph Theory · OpenGarden}
\subsection{Definitions and mathematical statement}
Let $G=(V,E)$ be a finite simple graph containing at least one edge. A
clique is a set of pairwise adjacent vertices, and $\omega(G)$ is the
largest size of such a set. A maximum clique has size $\omega(G)$; it is
not merely a clique maximal under inclusion. An odd hole is an induced
cycle of odd length at least five: its vertices have exactly the cycle
edges among them.

The conjecture asks whether every graph $G$ with no odd hole admits a
map $c:V\to\{0,1\}$ for which every maximum clique contains vertices of
both colours. Equivalently, it asks for a partition $V=X\sqcup Y$ with
\[
                 \omega(G[X])<\omega(G)
                 \quad\text{and}\quad
                 \omega(G[Y])<\omega(G).
\]
The colouring need not be proper, so edges within one colour class are
allowed. The forbidden configuration is specifically a monochromatic
clique of the largest size occurring anywhere in $G$.

The edge assumption makes explicit the source's word ``non-empty'': if
$G$ has vertices but no edges, each singleton is a maximum clique and
cannot use both colours. With $\omega(G)\ge2$ this degeneracy disappears.
When $\omega(G)=2$, the requirement becomes an ordinary proper
2-colouring, but for larger clique number it is weaker. Triangles are
allowed by the odd-hole hypothesis, and no prohibition on complements
of odd holes is imposed.
\subsection{Short English statement}
Can every odd-hole-free graph with an edge be coloured red and blue so that no largest clique is monochromatic?
\subsection{Research attempt}
\paragraph{Scope and process.}
This is a GPT-6 Astra Ultra research attempt dated 2 October 2026, using
primary-source browsing and elementary mathematical checks. The canonical
definition above is retained. Results below are restricted results or
reductions, not a claim of a new solution to the entire catalogue question.
No formal proof assistant or external mathematical referee checked this text.


\paragraph{Literature and scope.}
This is the Ho\`ang--McDiarmid two-division question. Chudnovsky and
Sivaraman [R1] prove divisibility results for restricted forbidden-induced-
subgraph classes, rather than for all odd-hole-free graphs. We prove
the target when clique and chromatic numbers agree, establish the
clique-number-two case directly, and give an explicit nonperfect family
where the required partition still exists. Maximum cliques, not all
inclusion-maximal cliques, are the objects being split throughout.

\paragraph{An exact certificate whenever $\chi=\omega$.}
Suppose $\chi(G)=\omega(G)=w\ge2$, and fix a proper colouring
with $w$ colours. Give all vertices of its first colour the new
colour red and give every other vertex blue. A clique of size $w$
uses $w$ distinct colours in the original proper colouring, hence
uses every one of them. It therefore contains exactly one red vertex
and $w-1$ blue vertices. This proves the desired conclusion whenever
$\chi=\omega$, without requiring the stronger assertion that every
induced subgraph has that equality.

In particular it applies to perfect graphs, but the hypothesis is not
forced merely by excluding odd holes. Complements of odd cycles of
length at least seven provide explicit counterexamples to that inference,
as analysed below. Thus a proof that simply calls every odd-hole-free
graph perfect would use a false premise.

\paragraph{The exact clique-number-two case.}
If $\omega(G)=2$, then $G$ has no triangle. If it had an odd
cycle, choose one of shortest odd length. Any chord would split it
into two smaller cycles, exactly one of them odd, contradicting the
choice. The cycle is consequently induced. Its length is at least
five because triangles are excluded, contradicting odd-hole-freeness.
So $G$ has no odd cycle and is bipartite: colour each connected
component by distance parity from a root; an edge within one parity
class would produce an odd closed walk and hence an odd cycle.
The bipartition splits every maximum clique, which in this case is
just an edge. The edge assumption in the canonical statement matters:
a singleton clique could not be split this way.

\paragraph{A nonperfect family handled by a different partition.}
Let $G$ be the complement of $C_{2m+1}$, with $m\ge3$ and the
vertices in cyclic order. An independent set in that cycle has size
at most $m$: mapping each selected vertex to its successor gives
an equally large disjoint set, so twice the size is at most $2m+1$.
Selecting alternate vertices attains $m$. Thus $\omega(G)=m$.
Every independent set in $G$ is a clique in the original cycle and
has size at most two, so $\chi(G)\ge m+1$. Pairing consecutive
cycle vertices and leaving one singleton gives $m+1$ colour classes,
proving equality. Hence this graph is not covered by $\chi=\omega$.

It has no odd hole, and in fact no induced cycle of any length at
least five. If $\ell$ selected vertices induced such a cycle in
$G$, each would have degree $\ell-3$ in the induced subgraph of
the original cycle. The original maximum degree is two, forcing
$\ell\le5$. At $\ell=5$, all five vertices would have original
degree two within the selected set, creating a five-cycle component
inside $C_{2m+1}$, which is impossible when $2m+1\ge7$.

Partition the cyclic order into consecutive blocks of lengths $m$
and $m+1$. In the original cycle each block induces a path, whose
largest independent set has size respectively $\lceil m/2\rceil$
and $\lceil(m+1)/2\rceil$. Both are at most $m-1$ for $m\ge3$.
The corresponding clique numbers in $G$ are therefore less than its
maximum $m$. Colouring the two blocks differently proves the target
for every graph in this family.

\paragraph{What a full solution would imply.}
The odd-hole-free class is hereditary. If the desired two-division
were available for every member, induction on clique number would
give $\chi(G)\le2^{\omega(G)-1}$ for every nonempty member: the
two parts each have smaller clique number, colour them separately,
and use disjoint palettes. The base case is an edgeless graph.
This deduction highlights the strength of a uniform partition rule;
it is not a derivation of that rule from a chromatic bound.

\paragraph{Verification and final status.}
The companion script [R2] checks clique numbers and the consecutive-block
partitions for complements of odd cycles of orders seven through
fifteen. The proofs cover all $m$ and distinguish the two different
meanings of maximal and maximum. \textbf{UNRESOLVED.} The missing
step is finding a suitable two-division for a general odd-hole-free
graph outside the explicitly certified families. No arbitrary proper
colouring provides that partition when more than $\omega(G)$ colours
are required.

\subsection{Sources}
\begin{itemize}
\item[S1] ULAM.AI and the UnsolvedMath Contributors, supplied problems.json, record 3251 (OPG-56230), OpenGarden collection. Catalogue identity and source transcription; CC BY 4.0.
\url{https://huggingface.co/datasets/ulamai/UnsolvedMath}.
\item[S2] Open Problem Garden, 2-colouring a graph without a monochromatic maximum clique. Original problem and discussion; the mathematical formulation below is expanded from the supplied source record.
\url{http://www.openproblemgarden.org/op/2_colouring_a_graph_without_a_monochromatic_maximum_clique}.

\item[R1] Maria Chudnovsky and Vaidy Sivaraman, \emph{Perfect divisibility
and 2-divisibility}, primary paper, abstract checked 2 October 2026.
\url{https://arxiv.org/abs/1704.06667}.
\item[R2] GPT-6 Astra Ultra, exact batch certificates, 2 October 2026.
Repository script \texttt{scripts/check\_attempt\_batch03\_colouring\_2026\_10\_02.py}.

\end{itemize}
\end{document}
